We summarize results for (i) the paired clean/OCR benchmark setup and (ii) baseline extraction performance on the OCR condition.

\subsection{Transcript-conditioned setup}
LongListBench now exposes two transcript conditions for the same rendered document: a clean canonical transcript and a noisy OCR transcript. This enables controlled clean-vs-OCR comparisons without changing the underlying layout, page structure, or ground truth. The paired setup is part of the release, but the quantitative baseline table in this revision focuses on the OCR condition, which matches the saved tiered reports currently included in the repository.

\subsection{Extraction regimes}
We compare three evaluation regimes that differ primarily in how they handle long OCR inputs. The first is full-context one-shot extraction, which submits the entire OCR transcript in one call and salvages complete incidents from truncated JSON when necessary. The second is auto-chunked extraction, which splits long OCR inputs into fixed token-budget chunks and extracts each chunk independently before concatenating the resulting incident lists. The third is agentic sandbox extraction, in which the model is given only the OCR inside a code-execution sandbox and decides for itself how to inspect and parse the document (typically by writing and running code) rather than following a fixed chunking policy. The first two regimes use Gemini 2.5; the agentic regime uses GPT-5.5. All regimes are evaluated on the same synchronized 80-document OCR snapshot from this repository.

\begin{table}[t]
\centering
\small
\caption{Comparison of the three extraction regimes on the synchronized 80-document OCR snapshot. Columns report corpus-level weighted field-value micro $F_1$ overall and on the extreme tier. The agentic regime uses GPT-5.5 while the one-shot and auto-chunked baselines use Gemini 2.5, so its gains reflect both regime and model.}
\label{tab:two-regimes}
\begin{tabular}{@{}llcc@{}}
\toprule
Regime & Model & Overall & Extreme \\
\midrule
One-shot & Gemini 2.5 & 27.4\% & 5.9\% \\
Auto-chunked & Gemini 2.5 & 84.8\% & 81.7\% \\
Agentic sandbox & GPT-5.5 & 90.6\% & 89.4\% \\
\bottomrule
\end{tabular}
\end{table}

Moving from full-context one-shot extraction to the local chunked baseline therefore improves overall weighted field-level $F_1$ by 57.4 points and extreme-tier weighted field-level $F_1$ by 75.8 points on the same snapshot. The one-shot regime remains strong on easy documents (97.2\%), but then degrades sharply to 74.6\% on medium, 44.4\% on hard, and 5.9\% on extreme. By contrast, the local chunked baseline reaches 97.3\% on easy, 96.5\% on medium, 87.7\% on hard, 71.0\% on detailed documents overall, and 95.9\% on table documents overall. In other words, chunking substantially reduces the severe long-document degradation seen in the one-shot regime, but the local chunked baseline still shows a pronounced gap between long detailed documents and long tables, with residual errors concentrated in record recovery, reading order, and layout-sensitive field completion. Because the strongest page-break and multi-column stressors are concentrated in detailed-format content in v1, this detailed-versus-table gap should not be interpreted as a pure format effect.

The agentic sandbox regime achieves the strongest results: 90.6\% overall weighted field-level $F_1$ (89.7\% recall, 91.4\% precision) with no extraction failures across all 80 documents. Critically, it largely closes the long-document gap, reaching 89.4\% on the extreme tier (500 incidents per document) versus 81.7\% for the chunked baseline and 5.9\% for one-shot, because the model navigates and parses the document programmatically rather than relying on a fixed chunking policy. By tier it reaches 97.0\% (easy), 95.2\% (medium), and 91.4\% (hard); by format it reaches 93.2\% on tables and 87.9\% on detailed documents, preserving the same detailed-versus-table ordering seen in the other regimes. Two caveats apply: the agentic regime uses a more capable model (GPT-5.5) than the Gemini baselines, so its margin reflects regime and model jointly rather than the regime in isolation; and it is more expensive, costing roughly \$0.40 per document (about \$33 for the full 80-document snapshot at current GPT-5.5 rates).

Qualitatively, errors often manifest as local field-level deviations (e.g., missing optional strings, numeric drift in financial breakdowns, or small identifier formatting mistakes) spread across an otherwise correct long list.

These findings suggest that the main open challenge for long-list extraction is robustly segmenting and populating full per-incident records under layout disruptions (page breaks, multi-column order, distractor tables, merged cells) and scale (hundreds of incidents). The paired clean/OCR benchmark design is intended to separate this difficulty from OCR-specific degradation in follow-up evaluations by comparing canonical and OCR conditions directly.
